\section{Introduction}
\label{sec:intro}

% Paragraph 1 — Problem & motivation
Online book platforms face a dual challenge: users who know what they want
need precise retrieval, while passive browsers benefit from proactive
suggestions.
Most systems address only one of these modes, leaving an
\emph{interaction-content gap}—behavioral signals (co-purchase patterns) are
decoupled from semantic content understanding (title, description, cover
image).

% Paragraph 2 — Limitations of prior work
Collaborative filtering methods excel at capturing collective behavior but
suffer from cold-start and fail to leverage rich multimodal content~\cite{he2020lightgcn}.
Content-based methods using text or image encoders lack the relational
structure that behavioral co-occurrence provides.
Sequential models~\cite{kang2018self} capture preference evolution but typically ignore
visual modality.

% Paragraph 3 — Our approach & contributions
We propose a \textbf{Dual-Mode Multimodal Recommendation System} that bridges
these gaps through:
\begin{itemize}
  \item \textbf{Mode~1 (Active Search):} multilingual query-driven retrieval
        (19 languages via NLLB-200) fusing BGE-M3 text embeddings and CLIP
        image embeddings via Reciprocal Rank Fusion (RRF), with optional BM25
        keyword boosting and multi-layer LRU caching for sub-25\,ms repeat
        query latency.
  \item \textbf{Mode~2 (Passive Recommendation):} two complementary pipelines—
        (a)~Cleora graph embeddings for behavioral similarity
        (\emph{People Also Buy}), and (b)~DIF-SASRec sequential personalization
        (\emph{You Might Like}).
  \item A \textbf{shared user profile} updated after every interaction,
        bridging both modes with temporal weighting.
\end{itemize}

% Paragraph 4 — Paper structure
The remainder of this paper is organized as follows.
Section~\ref{sec:related} reviews related work.
Section~\ref{sec:overview} describes the system architecture.
Section~\ref{sec:method} details each component.
Sections~\ref{sec:exp} and~\ref{sec:results} present experimental setup and
results.
Section~\ref{sec:conclusion} concludes.
